% P10, 15.09.2026: konkreter Browserweg; LocalReviewServerHost.cs und ReviewUiFlow.cs.
% M10, 11.09.2026: Leserführung, Kapiteltrennung und belegte Reichweite.
% PipelineFullRunner.EventPump.cs:80--212 (Dispatch und gate-spezifische Regeln);
% BacklogStageExecutors.cs:27--40; PipelineFullRunner.GateResponder.cs; HitlShell.cs.
% B-32: zentraler Dispatch ist keine globale Replay-Regel. Ingest-/UI-Dateivertrag in 6.6.

\section{Human Gates mit Request Ports}
\label{sec:realisierung-human-gates}

Die menschlichen Entscheidungspunkte aus
Abschnitt~\ref{sec:governance-mutation} verwenden \emph{Request Ports}.
Über diese MAF-Schnittstellen gibt ein Workflow eine typisierte Anfrage
nach außen und wartet auf eine dazu passende Antwort. Am Beispiel der
Einordnungsstufe aus
Abschnitt~\ref{sec:realisierung-agent-executor-workflow} wird daraus
ein konkreter Übergang: Nach erfolgreicher Prüfung sendet die
Finalisierung den Operationsplan an den Gate-Port. Dessen Anfrage
erscheint im Ereignisstrom des Laufs. Eine zentrale Ereignisschleife
ordnet sie anhand der Port-Kennung dem zuständigen Antwortverfahren zu.

Im hier beschriebenen Betrieb mit persistierter Gate-Pause wird
zunächst geprüft, ob bereits eine eingereichte Antwort vorliegt.
Andernfalls sichert der Lauf einen Checkpoint mit der offenen Anfrage
und hinterlegt einen \emph{Lauf-Zeiger}, der Sitzung, Checkpoint und
Gate identifiziert. Damit kann der ausführende Prozess beendet
werden, ohne den vorgesehenen Fortsetzungspunkt zu verlieren.

Für die Review-Oberfläche startet der jeweilige Review-Aufruf einen
lokalen Webserver, der die Bearbeitung im Browser ermöglicht. Je nach
Gate zeigt die Oberfläche die vorgelegten Vorschläge, zugehörige Belege
oder Bestandsinformationen und die verfügbaren Entscheidungsoptionen.
Die Eingaben des Menschen werden als Entscheidungen gespeichert.
Alternativ unterstützt der Steward die Einreichung im Dialog, wie
Abschnitt~\ref{sec:realisierung-steward} erläutert.

Bei der Wiederaufnahme wird die gespeicherte Entscheidung geladen und
in die zum Gate passende Antwort übertragen: je nach Gegenstand etwa
eine Auswahl von Elementen, eine überarbeitete Fassung oder eine
Konfliktentscheidung. Erst die Antwort ermöglicht die nachgelagerte
Verarbeitung.

Die deklarierten Experiment-Betriebsarten beantworten Anfragen dagegen
nach gate-spezifischen Regeln. Die zentrale Verteilung bedeutet hier
keine einheitliche Entscheidungsregel: Beispielsweise weist das
Cluster-Review beim Einspielen gespeicherter Entscheidungen
(\emph{Replay}) unbekannte Elemente ab; im Backlog-Review gilt ein
fehlender Eintrag dagegen als Annahme ohne Bearbeitung.
Architekturklassifikation und ADR-Freigabe übernehmen in diesem Modus die
Vorschläge. Replay ist somit kein globaler Nachweis unveränderter
menschlicher Einzelentscheidungen. Offene fachliche Konflikte werden
in den Experimentmodi ausdrücklich vertagt. Leere Entscheidungsmengen
umgehen die Anfrage über einen protokollierten Leerfall.

Das Human Gate ist von Werkzeugzustimmung und Invariantenprüfung
abzugrenzen: Es erfasst die fachliche Entscheidung über den vorgelegten
Gegenstand. Die Werkzeugzustimmung erlaubt eine einzelne Handlung
(Abschnitt~\ref{sec:realisierung-tools-mcp}); die Speicherprüfung
kontrolliert die strukturelle Zulässigkeit der resultierenden Mutation
(Abschnitt~\ref{sec:realisierung-persistenz-aussenkopplung}). Die
Gate-Übersicht steht in Anhang~\ref{anh:gate-matrix}. Wie der Lauf über
die beschriebene Pause hinweg erhalten bleibt, zeigt der nächste
Abschnitt.
